% -*- coding: utf-8 -*-
% ============================================================
%  main.tex —— ZJU Beamer 模板功能演示
%  编译：latexmk -xelatex main.tex（需 XeLaTeX，编译两遍以更新页码/目录）
%  改内容：直接修改本文件；样式与组件见 zju_beamer_pro.sty
% ============================================================
\PassOptionsToPackage{quiet}{fontspec}
\PassOptionsToPackage{dvipsnames}{xcolor}
\documentclass[10pt,aspectratio=169]{beamer}
\usepackage{zju_beamer_pro}
\graphicspath{{./figures/}}

% ---------------- 首页信息（改成你自己的） ----------------
\title{浙江大学 Beamer 模板}
\subtitle{ZJU Beamer Pro · 简洁学术风演示模板}
\author{张三\\[1mm]{\footnotesize 计算机科学与技术专业}}
\institute{浙江大学计算机科学与技术学院}
\date{\today}

\begin{document}

% ============================ 封面 ============================
\begin{frame}[plain]
  \titlepage
\end{frame}

% ============================ 目录 ============================
\begin{frame}{目录}{CONTENTS}
  \vspace{2mm}
  \begin{minipage}[t]{0.48\textwidth}
    \vspace{0pt}%
    \tableofcontents[sections={1-2}]
  \end{minipage}\hfill
  \begin{minipage}[t]{0.48\textwidth}
    \vspace{0pt}%
    \tableofcontents[sections={3-4}]
  \end{minipage}
\end{frame}

% ============================================================
%   01 快速上手
% ============================================================
\section{快速上手}

\begin{frame}[fragile]{快速上手}{QUICK START}
  \begin{minipage}[t]{0.48\textwidth}
    \vspace{0pt}%
    \begin{zjubox}{blue}{环境依赖}
      {\footnotesize
      \begin{denseitem}
        \item TeX Live 2023+（Overleaf 亦可）
        \item 编译器：\textbf{XeLaTeX}
        \item 字体已随 \texttt{fonts/} 分发，无需安装
        \item 西文 Fira Sans 来自 TeX Live 的 \texttt{fira} 宏包
      \end{denseitem}}
    \end{zjubox}
    \begin{zjubox}{gray}{Overleaf 使用}
      {\footnotesize 新建项目 $\rightarrow$ 上传 Release 中的模板压缩包
      $\rightarrow$ 菜单中将编译器切换为 \textbf{XeLaTeX}
      $\rightarrow$ 重新编译两次。}
    \end{zjubox}
  \end{minipage}\hfill
  \begin{minipage}[t]{0.48\textwidth}
    \vspace{0pt}%
    \begin{zjubox}{green}{Overleaf 之外}
      {\footnotesize 也可以在 Overleaf 新建空白项目后上传整个仓库，
      菜单中把编译器切换为 \textbf{XeLaTeX} 后编译两次。}
    \end{zjubox}
    \begin{zjubox}{orange}{在哪里改内容}
      {\footnotesize 所有正文都在 \texttt{main.tex}：
      封面信息在文件开头的 \texttt{\string\title} 等命令中，
      正文按 section 组织，组件用法见本演示各页源码。}
    \end{zjubox}
  \end{minipage}
  \vspace{0.6mm}
  \begin{zjucode}{bash}{本地编译（编译两遍）}
latexmk -xelatex main.tex
latexmk -xelatex main.tex
  \end{zjucode}
  \begin{takeaway}
    封面、章节过渡页、致谢页由模板自动生成，只需填好首页信息即可。
  \end{takeaway}
\end{frame}

% ============================================================
%   02 版式与布局
% ============================================================
\section{版式与布局}

\begin{frame}{版式样式}{LAYOUT STYLES}
  \vspace{1mm}
  \begin{minipage}[t]{0.48\textwidth}
    \vspace{0pt}%
    \begin{zjubox}{blue}{自动生成的版式}
      {\small
      \begin{denseitem}
        \item \textbf{封面}：求是蓝校徽校名 + 浅蓝色带大标题
        \item \textbf{章节过渡页}：\texttt{\string\section} 自动插入
        \item \textbf{致谢页}：\texttt{\string\zjuthanks} 一行调用
        \item \textbf{页眉横线}：标题 + 渐变细线 + 渐变校名
      \end{denseitem}}
    \end{zjubox}
  \end{minipage}\hfill
  \begin{minipage}[t]{0.48\textwidth}
    \vspace{0pt}%
    \begin{zjubox}{green}{页标题与副标题}
      {\small 每页标题带副标题参数：}

      \vspace{0.8mm}
      {\small\color{black!60}
      \texttt{\string\begin\{frame\}\{标题\}\{SUBTITLE\}}}

      \vspace{0.8mm}
      {\small 副标题以灰色小字显示在标题右侧，
      建议使用英文大写缩写。}
    \end{zjubox}
  \end{minipage}
  \vspace{1.2mm}
  \begin{zjubox}{gray}{整页纯白背景}
    {\small 内容页无任何背景装饰，适合高密度学术内容；
    全出血封面与章节页自带留白设计，16:9 版面。}
  \end{zjubox}
  \begin{takeaway}
    本页右侧即「标题 + 渐变横线 + 渐变校名」的页眉样式。
  \end{takeaway}
\end{frame}

\begin{frame}{多栏布局}{MULTI-COLUMN}
  % ---- 两栏 ----
  \begin{minipage}[t]{0.475\textwidth}
    \vspace{0pt}%
    \begin{zjubox}{blue}{两栏：minipage + hfill}
      {\footnotesize 用 \texttt{minipage[t]} 并排内容，
      首行加 \texttt{\string\vspace\{0pt\}} 保证顶部对齐；
      栏间用 \texttt{\string\hfill} 撑开。}
    \end{zjubox}
  \end{minipage}\hfill
  \begin{minipage}[t]{0.475\textwidth}
    \vspace{0pt}%
    \begin{zjubox}{sky}{任意分栏比例}
      {\footnotesize 栏宽自由组合：
      0.475 + 0.475、0.3 + 0.3 + 0.3、
      0.6 + 0.38 等均可，
      组合 \texttt{zjubox}、\texttt{imgcard}、图表皆可。}
    \end{zjubox}
  \end{minipage}
  \vspace{1.2mm}

  % ---- 三栏 ----
  \begin{minipage}[t]{0.31\textwidth}
    \vspace{0pt}%
    \begin{zjubox}{orange}{三栏之一}
      {\footnotesize 三栏时建议 0.31 + 0.31 + 0.31。}
    \end{zjubox}
  \end{minipage}\hfill
  \begin{minipage}[t]{0.31\textwidth}
    \vspace{0pt}%
    \begin{zjubox}{purple}{三栏之二}
      {\footnotesize 各栏首行同样加
      \texttt{\string\vspace\{0pt\}}。}
    \end{zjubox}
  \end{minipage}\hfill
  \begin{minipage}[t]{0.31\textwidth}
    \vspace{0pt}%
    \begin{zjubox}{steel}{三栏之三}
      {\footnotesize 高度不同时顶部自动对齐。}
    \end{zjubox}
  \end{minipage}

  \begin{takeaway}
    切记：\texttt{minipage[t]} 内第一个元素前加 \texttt{\string\vspace\{0pt\}}，否则顶部不对齐。
  \end{takeaway}
\end{frame}

% ============================================================
%   03 预定义组件与图片
% ============================================================
\section{组件与图片}

\begin{frame}{功能框与要点条}{COLORED BOXES}
  \begin{zjubox}{blue}{zjubox：全色系功能框}
    {\small 七种色键全部映射为蓝色阶梯：
    \badge[zjublue]{blue}\ \badge[zjubluelight]{orange}\ \badge[zjusky]{green}\ \badge[zjuindigo]{purple}\ \badge[zjuindigo]{red}\ \badge[zjusky]{gold}\ \badge[zjusteel]{gray}，
    用法 \texttt{\string\begin\{zjubox\}\{色键\}\{标题\}}。
    标题条 + 嫩蓝浅底 + 圆角无边框，蓝色系视觉统一。}
  \end{zjubox}
  \vspace{1mm}
  \begin{minipage}[t]{0.475\textwidth}
    \vspace{0pt}%
    \begin{zjubox}{purple}{配合徽章 \string\badge}
      {\small 徽章用于标签化信息：\badge[zjublue]{标签一}\;
      \badge[zjubluelight]{标签二}\; \badge[zjuindigo]{重点}}
    \end{zjubox}
  \end{minipage}\hfill
  \begin{minipage}[t]{0.475\textwidth}
    \vspace{0pt}%
    \begin{zjubox}{steel}{高亮命令 \string\hl}
      {\small 正文强调用 \hl{蓝色加粗}；
      takeaway 等深色底上会自动切换为浅色。}
    \end{zjubox}
  \end{minipage}
  \begin{takeaway}
    takeaway：全宽渐变要点条，用于每页收束一句话结论。
  \end{takeaway}
\end{frame}

\begin{frame}{数据卡片与时间线}{STATCARDS \& TIMELINE}
  % ---- 数据卡片 ----
  \zjustatwidth=4.55cm
  \begin{center}
    \statcard[blue]{98.6}{示例百分制成绩}\hfill
    \statcard[orange]{4.5\,/\,5.0}{示例均绩}\hfill
    \statcard[green]{12 项}{示例条目计数}
  \end{center}
  \vspace{0.4mm}
  % ---- 时间线 + 奖项卡片 ----
  \begin{minipage}[t]{0.475\textwidth}
    \vspace{0pt}%
    \begin{zjubox}{blue}{zjutimeline 时间线}
      \begin{zjutimeline}
        \titem{gray}{2024}{阶段一}{示例条目}
        \titem{silver}{2025}{阶段二}{示例条目}
        \titem{gold}{2026}{阶段三}{示例条目}
      \end{zjutimeline}
    \end{zjubox}
  \end{minipage}\hfill
  \begin{minipage}[t]{0.475\textwidth}
    \vspace{0pt}%
    \begin{zjubox}{orange}{awardcard 奖项卡片}
      \awardcard[gold]{trophy}{一等奖}{示例竞赛名称}
      {示例描述 · 团队成员}
    \end{zjubox}
  \end{minipage}
  \begin{takeaway}
    statcard 可用 \string\hfill 并排，宽度由 \string\zjustatwidth 统一控制。
  \end{takeaway}
\end{frame}

% ============================================================
%   04 代码与公式
% ============================================================
\section{代码与公式}

\begin{frame}[fragile]{代码样式}{CODE}
  \begin{zjucode}{Python}{Python 示例}
def fibonacci(n: int) -> int:
    a, b = 0, 1
    for _ in range(n):
        a, b = b, a + b
    return a
# keyword / comment / string 均有独立配色
  \end{zjucode}
  \vspace{0.4mm}
  \begin{minipage}[t]{0.475\textwidth}
    \vspace{0pt}%
    \begin{zjubox}{green}{样式说明}
      {\footnotesize
      \begin{denseitem}
        \item 基于 listings 宏包，浅冰蓝底 + 细边框
        \item 关键字深蓝加粗，行号可开关
      \end{denseitem}}
    \end{zjubox}
  \end{minipage}\hfill
  \begin{minipage}[t]{0.475\textwidth}
    \vspace{0pt}%
    \begin{zjubox}{steel}{使用提示}
      {\footnotesize 用模板提供的 \texttt{zjucode} 代码框，
      无需 \texttt{[fragile]}；也可直接使用
      \texttt{lstlisting}（此时页面需 \texttt{[fragile]}）。}
    \end{zjubox}
  \end{minipage}
  \begin{takeaway}
    listings 已预设配色，直接使用 \texttt{lstlisting} 环境即可。
  \end{takeaway}
\end{frame}

\begin{frame}{公式样式}{MATHEMATICS}
  均值不等式与常用推导：
  \begin{block}{常用不等式}
    \vspace{-1.5mm}
    \begin{align}
      \frac{a+b}{2} &\geq \sqrt{ab},
      & \forall a,b>0 \\
      \sum_{i=1}^{n} i^2 &= \frac{n(n+1)(2n+1)}{6},
      & n \in \mathbb{N}
    \end{align}
  \end{block}
  \vspace{0.5mm}
  行内公式保持正文节奏：梯度 $\nabla_\theta \mathcal{L}(\theta)$、
  注意力 $\mathrm{softmax}(QK^{\top}/\sqrt{d})V$。
  \begin{takeaway}
    数学字体为衬线（\texttt{onlymath}），与无衬线正文形成对比。
  \end{takeaway}
\end{frame}

\begin{frame}{多栏图片}{FIGURES}
  \begin{minipage}[t]{0.475\textwidth}
    \vspace{0pt}%
    \centering
    \imgcard[6.9cm]{demo-a.png}{示例图 A：\texttt{imgcard} 自带边框与图注}
  \end{minipage}\hfill
  \begin{minipage}[t]{0.475\textwidth}
    \vspace{0pt}%
    \centering
    \imgcard[6.9cm]{demo-b.png}{示例图 B：宽度建议 \texttt{\string\linewidth}}
  \end{minipage}
  \vspace{1.2mm}
  \begin{zjubox}{gray}{图文混排}
    {\small 图片栏与文字栏组合：
    左侧 \texttt{imgcard}，右侧 \texttt{zjubox} 或列表；
    \texttt{imgcard} 的第一个参数控制内容宽度。}
  \end{zjubox}
  \begin{takeaway}
    双图并排：两个 \texttt{minipage[t]} 各放一个 \texttt{imgcard}，中间 \string\hfill。
  \end{takeaway}
\end{frame}

% ============================ 致谢 ============================
\begin{frame}[plain]
  \zjuthanks
\end{frame}

\end{document}
